% ============================================================
%  第2章：单纯形法
% ============================================================
\section{第2章 \quad 单纯形法}

\subsection{2.1 从枚举到智能搜索}

第1章的核心结论：标准型LP的最优解一定在某个基本可行解（BFS，即可行域的极点）上。理论上，只需枚举全部 $\binom{n}{m}$ 个BFS，比较目标值挑最优即可。但这不现实——$\binom{50}{25}\approx 10^{14}$，宇宙年龄都不够用。

\textbf{单纯形法（Simplex Method）是一种"聪明地遍历"BFS的方法}：从一个初始BFS出发，每一步只考虑\emph{能让目标值下降}的相邻BFS，沿着可行域的棱走过去。绝大多数候选BFS根本不用看。

\textbf{核心思想}：在顶点间沿棱跳跃，每一步都让目标值\emph{严格下降}（对$\min$问题）。因为BFS总数有限且目标严格单调，算法必然在有限步终止。几何上，可行域是一个凸多面体，单纯形法从某个顶点出发，沿着边界棱的方向移动到相邻顶点，目标值不断改善。（课件 p5.1）

\subsection{2.2 检验数：衡量"进基是否划算"}

\subsubsection{从基矩阵分解出发}

设当前基为 $B$（$m\times m$ 可逆），非基矩阵为 $N$。将矩阵分块：
\[
A = [B \mid N],\quad x = \binom{x_B}{x_N},\quad c = \binom{c_B}{c_N}
\]
约束 $Ax=b$ 写为 $Bx_B + Nx_N = b$。因为 $B$ 可逆，基变量可由非基变量表出：
\[
x_B = B^{-1}b - B^{-1}N x_N
\]
代入目标函数：
\[
\begin{aligned}
z &= c_B^T x_B + c_N^T x_N
   = c_B^T(B^{-1}b - B^{-1}N x_N) + c_N^T x_N \\
  &= c_B^T B^{-1}b + (c_N^T - c_B^T B^{-1}N)\,x_N
\end{aligned}
\]
令 $z_0 = c_B^T B^{-1}b$（当前目标值），则：
\[
z = z_0 + \sum_{j \in \text{非基}} (c_j - c_B^T B^{-1} a_j)\,x_j
\]

\textbf{定义（检验数 / Reduced Cost）}

非基变量 $x_j$ 的\emph{检验数}定义为：
\[
\sigma_j = c_B^T B^{-1} a_j - c_j
\]
其中 $z_j = c_B^T B^{-1} a_j$ 是"引入一单位 $x_j$ 的机会成本"——把 $x_j$ 从0提升到1，为了保持等式成立，基变量必须调整；$z_j$ 就是这些调整带来的目标变化。

\textbf{最优性条件}（对 $\min$ 问题）：若所有非基变量的检验数 $\sigma_j \le 0$，则当前BFS为最优解。

\textbf{无界条件}：若存在 $\sigma_k > 0$ 但对应列 $B^{-1}a_k \le 0$（全为非正），则 $x_k$ 可无限增大使目标无限下降，问题无下界（无有限最优解）。

\textbf{检验数的直觉：} $\sigma_j > 0$ 意味着增大 $x_j$（让非基变量进基）能使目标\emph{下降}（对 $\min$ 问题，$\sigma_j>0$ 说明 $x_j$ 有改善潜力）。$\sigma_j$ 就是"让 $x_j$ 增加1单位的\emph{净收益}"——$z_j$（机会成本减省）减去 $c_j$（新增成本）。$z = z_0 - \sum \sigma_j x_j$，$\sigma_j$ 愈大，$x_j$ 每增1单位目标下降愈多。（课件 p5.6）

\subsection{2.3 单纯形表：手工计算的载体}

单纯形表把基矩阵信息、检验数、当前解值压缩在一张表格里，便于手工迭代。

\textbf{单纯形表的结构}

\[
\begin{array}{c|c|cccccc|c}
c_B & \text{基} & x_1 & x_2 & \cdots & x_j & \cdots & x_n & b \\ \toprule
c_{B_1} & x_{B_1} & \bar{a}_{11} & \bar{a}_{12} & \cdots & \bar{a}_{1j} & \cdots & \bar{a}_{1n} & \bar{b}_1 \\
c_{B_2} & x_{B_2} & \bar{a}_{21} & \bar{a}_{22} & \cdots & \bar{a}_{2j} & \cdots & \bar{a}_{2n} & \bar{b}_2 \\
\vdots  & \vdots  & \vdots  & \ddots & \vdots  & \ddots & \vdots  & \vdots  \\
c_{B_m} & x_{B_m} & \bar{a}_{m1} & \bar{a}_{m2} & \cdots & \bar{a}_{mj} & \cdots & \bar{a}_{mn} & \bar{b}_m \\ \midrule
- & \sigma_j & \sigma_1 & \sigma_2 & \cdots & \sigma_j & \cdots & \sigma_n & z_0
\end{array}
\]
其中 $\bar{A}=B^{-1}A$，$\bar{b}=B^{-1}b$。末行检验数，末格为 $z_0$（当前目标函数值）。$c_B$ 列在最优表中给出对偶变量（影子价格）。

\textbf{读出信息的规则}：
\begin{itemize}
  \item $c_B$ 列（最左列）= 每个基变量在目标函数中的系数——是计算检验数和 $z_0$ 的关键输入
  \item 基变量对应的列是单位向量（如基变量 $x_4$ 列的对应行是1，其余行0），其检验数恒为0
  \item $b$ 列 = 当前基变量的取值（非基变量均为0）
  \item 末行末格 = $z_0$（即当前目标函数值，$z_0 = c_B \cdot b$）
\end{itemize}
（课件 p5.10）

\subsubsection{检验数的表中计算公式}

前面定义 $\sigma_j = c_B^T B^{-1} a_j - c_j$ 用了矩阵语言。在单纯形表中，检验数完全用表中数字计算：

\textbf{核心公式（表中版本）：}
\[
\sigma_j = \sum_{i=1}^{m} c_{B_i} \!\cdot\! \bar{a}_{ij} \;-\; c_j
\]
其中 $c_{B_i}$ 是第 $i$ 行基变量的目标系数，$\bar{a}_{ij}$ 是表中第 $i$ 行第 $j$ 列的系数，$c_j$ 是 $x_j$ 的原目标系数（$\min$ 标准型）。

\textbf{一句话记住：} 把基变量的"价格"按对应的行系数加权求和，就是"机会成本"$z_j$；再减掉这个变量自己的价格 $c_j$，就是检验数。

\textbf{检验数计算口诀：}
\begin{enumerate}
  \item 找到表左边\emph{基变量列}，写下它们的目标系数 $c_B$
  \item 对每一列 $j$：\textbf{同行相乘再求和，减掉原系数} $\; \sigma_j = \sum_{i} c_{B_i}\bar{a}_{ij} - c_j$
  \item 末格同理：$z_0 = \sum_{i} c_{B_i} \bar{b}_i$
  \item 基变量列必为 $0$（单位向量 + 公式保证），算出来不放心时可以拿来验算
\end{enumerate}

这个方法对任何一张单纯形表都成立——初始表、迭代中间表、最优表、含 $M$ 的表，只要是正确的单纯形表，公式一模一样。

\subsubsection{比值的计算：谁先被逼到零}

确定了进基变量 $x_k$ 之后，下一步要决定\emph{哪个基变量离基}。这就是"比值"的作用。

\textbf{核心公式：}
\[
\theta_i = \frac{\bar{b}_i}{\bar{a}_{ik}}, \quad \text{只对 } \bar{a}_{ik} > 0 \text{ 的行计算}
\]
取所有 $\theta_i$ 中\emph{最小的}，对应的第 $i$ 行基变量离基。

\textbf{为什么是这个公式？} 基本可行解要求所有基变量 $\ge 0$。当进基变量 $x_k$ 从 $0$ 增加到 $\Delta$ 时，第 $i$ 行基变量变为：
\[
x_{B_i} = \bar{b}_i - \bar{a}_{ik} \cdot \Delta
\]
要让 $x_{B_i} \ge 0$，必须有 $\Delta \le \bar{b}_i / \bar{a}_{ik}$（仅当 $\bar{a}_{ik} > 0$ 时构成上限约束）。取所有上限中最小的——第一个被逼到零的基变量就是"瓶颈"，它离基。

\textbf{比值计算口诀：}
\begin{enumerate}
  \item 锁定进基列 $x_k$，从上到下扫每一行
  \item 系数 $\bar{a}_{ik} \le 0$？跳过——这行不限制 $x_k$ 的增长
  \item 系数 $\bar{a}_{ik} > 0$？计算 $\theta_i = \bar{b}_i / \bar{a}_{ik}$
  \item 取最小的 $\theta_i$——对应的基变量离基，该行该列的交点就是\emph{主元}
  \item 若全列无正系数（所有 $\bar{a}_{ik} \le 0$），则 $x_k$ 可以无限增大，问题\emph{无下界}（无有限最优解），算法终止
\end{enumerate}

比值和检验数是单纯形法的两个核心计算——检验数决定"谁进"，比值决定"谁出"。两者都只需要表中已有的数字，不需要额外推导。

\subsection{2.4 一次单纯形迭代：三步走透}

假设当前BFS非退化，一次完整的单纯形迭代由三步组成。

\textbf{方法（单纯形法一次迭代）}

\begin{enumerate}
  \item \textbf{进基选择}：在检验数行中，找\emph{最大的正检验数} $\sigma_k$，对应的非基变量 $x_k$ 进基。若所有检验数 $\le 0$，当前解已最优，停止。
  \item \textbf{离基选择（最小比值原则）}：对 $x_k$ 列中所有 $\bar{a}_{ik} > 0$ 的行，计算 $\theta_i = \bar{b}_i / \bar{a}_{ik}$，取最小 $\theta_i$ 对应的基变量 $x_r$ 离基。若该列无正系数，问题无下界（目标可无限下降），停止。
  \item \textbf{转轴运算（Pivot）}：以 $\bar{a}_{rk}$ 为主元，对单纯形表进行行变换——主元行除以 $\bar{a}_{rk}$ 使主元变为1，然后用此行消去其他所有行（包括检验数行）的 $x_k$ 列，使该列变成第 $r$ 个分量为1的单位向量。
\end{enumerate}

下面逐条展开因果。

\subsubsection{进基：为什么选最大正检验数？}

\textbf{从哪来？} 检验数 $\sigma_j$ 衡量的是"让 $x_j$ 增加1单位，目标能改善多少"。对 $\min$ 问题，$\sigma_j>0$ 意味着增加 $x_j$（从0变正）能让目标下降。

\textbf{为什么做？} 选最大正检验数就是选"单位增量收益最大"的变量进基——贪心思想，每步选下降最快的方向。实际单纯形法也常用其他规则（如选第一个正检验数），理论效率相近。

\textbf{得到什么？} 确定进基变量 $x_k$，即下一步要增大哪个非基变量。

\textbf{例：} 设检验数行为 $(\sigma_1,\sigma_2,\sigma_3,\sigma_4,\sigma_5) = (3,-1,0,0,5)$。最大正检验数是 $\sigma_5=5$，所以 $x_5$ 进基。这意味着让 $x_5$ 从0增大，每增1单位目标下降5单位——最划算。

\subsubsection{离基：为什么取最小正比值？}

\textbf{从哪来？} 当 $x_k$ 增大时，基变量随之变化：$x_{B_i} = \bar{b}_i - \bar{a}_{ik} \cdot \Delta$。如果 $\bar{a}_{ik}>0$，$x_{B_i}$ 随 $x_k$ 增大而减少。

\textbf{为什么做？} 基变量必须 $\ge 0$（可行性要求）。$x_k$ 最多能增到多大？对每个 $\bar{a}_{ik}>0$ 的行，上限是 $\bar{b}_i/\bar{a}_{ik}$（再大该基变量就变负了）。取所有上限中的最小值——最先被逼到0的基变量就是"瓶颈"，它离基。

\textbf{得到什么？} 确定离基变量 $x_r$ 和主元 $\bar{a}_{rk}$。$x_k$ 增大到 $\theta=\min_i\{\bar{b}_i/\bar{a}_{ik}\mid \bar{a}_{ik}>0\}$，$x_r$ 降为0出基。

\textbf{为何只看正系数？} $\bar{a}_{ik}\le 0$ 时，$x_k$ 增大反而让 $x_{B_i}$ 增大或不变——这条约束不限制 $x_k$ 的增长，不是瓶颈。

\textbf{例：} 设 $x_k=x_3$，其列系数为 $(2, -1, \frac{1}{2})^T$，右端 $b=(6,4,3)^T$。比值：$6/2=3$，$4/(-1)$ 跳过（负数），$3/(1/2)=6$。最小是 $3$，对应第一行，该行基变量 $x_1$ 离基，主元为 $\bar{a}_{13}=2$。

\subsubsection{转轴：为什么这样做？}

\textbf{从哪来？} 已有进基变量 $x_k$ 和离基变量 $x_r$，主元 $\bar{a}_{rk}$。

\textbf{为什么做？} 转轴运算是一次基变换的代数实现——用进基列替换离基列，使新基仍然是单位阵。主元行除以主元让新基变量列在对应位置为1；用此行消去其他行让该列其他位置为0。

\textbf{得到什么？} 新单纯形表，新BFS，目标值下降（或不变，退化时）。检验数行也同时更新。

\textbf{为什么检验数行也能用行变换更新？（而不是每次重算）}

这是初学者最容易困惑的地方。答案分两步：

\textbf{第一步：检验数行本质上就是一个线性方程。}

把目标函数改写为 $z + \sum_{j}\sigma_j x_j = z_0$，它就是关于变量 $(z,x_1,\dots,x_n)$ 的一个线性方程——和约束方程 $x_{B_i} + \sum_j \bar{a}_{ij}x_j = \bar{b}_i$ 地位完全相同。整个单纯形表就是一个 $(m+1)$ 行的线性方程组。

\textbf{第二步：线性方程组的行变换不改变解集。}

对方程组做"某行减去另一行的倍数"（Gauss消元），等价的方程组有相同的解。所以当我们用主元行消去第 $k$ 列其他行时，\emph{检验数行也参与同样的消元操作}——这保证新表仍然是等价的方程组，其中的 $\sigma_j'$ 就是新基下的检验数，$z_0'$ 就是新目标值。

行变换不是"偷懒"，而是 Gauss 消元的严格推论——它等价于对新基 $B'$ 重新计算 $B'^{-1}A$ 和 $c_{B'}^T B'^{-1}A - c$，只是省去了矩阵求逆的麻烦。

（课件 p5.4, p5.5）

\subsection{2.5 初始基本可行解从哪里来？}

单纯形法需要一个初始BFS才能启动。但并非所有标准型LP的系数矩阵都包含现成的单位子阵。

\textbf{什么时候直接有？} 全部是 $\le$ 约束且 $b\ge 0$：每个约束加的松弛变量系数为 $+1$，天然构成 $m$ 阶单位阵。松弛变量就是初始基变量，初始BFS = $(0,\dots,0,b_1,\dots,b_m)^T$。

\textbf{什么时候没有？} 两种情况导致缺基变量列：
\begin{itemize}
  \item 约束中有 $\ge$：剩余变量系数是 $-1$，不能当单位向量
  \item 约束中有 $=$：根本没有自然添加的变量
\end{itemize}
这两种情况下，系数矩阵缺少足够的 $+1$ 单位列，无法直接读出初始BFS。这就需要用大M法或两阶段法来"造"初始基。（课件 p5.8）

\subsection{2.6 大M法}

\subsubsection{问题根源：为什么需要人工变量？}

某个约束缺少"系数为$+1$且在其他行系数为0"的变量。例如 $\ge$ 约束减去剩余变量后，剩余变量的系数是 $-1$，把它放进基意味着 $x_B = -b < 0$（因为 $b>0$），不是可行解。

\textbf{人工变量的思路：} 既然缺 $+1$ 列，我们就\emph{硬塞一个 $+1$ 列进去}。在每个缺基变量的行引入一个\emph{人工变量} $x_a\ge 0$，系数正好是 $+1$。这样人工变量加上原有的松弛变量，恰好凑出完整的 $m$ 阶单位阵——初始基就有了。

\subsubsection{核心矛盾：人工变量是"假的"}

人工变量不是原问题的变量——它只是我们为了凑初始基\emph{人为添加}的。在真正的最优解中，人工变量必须等于0（否则原约束被篡改了）。但单纯形法不知道谁是"假的"——如果不加干预，算法会让 $x_a>0$ 留在解中。

\subsubsection{大M法的解决策略：惩罚驱赶}

\textbf{为什么加M？} 在目标函数中对每个人工变量加上一个\emph{巨大的惩罚系数 $M>0$}（对 $\min$ 问题加 $+M x_a$，对 $\max$ 问题加 $-M x_a$）。$M$ 是"非常大"的正数——大到任何不含人工变量的可行解目标值，都优于含正人工变量的可行解。

\textbf{M为什么能赶走人工变量？} 考虑 $\min$ 问题。目标变为 $\min\; c^T x + M\sum x_a$。如果某人工变量 $x_a>0$，目标中 $M x_a$ 这一项就极大，整个目标值被拖累得很差。单纯形法要最小化目标——它会主动让所有正检验数对应的人工变量尽快离基，因为 $M$ 的系数巨大，人工变量的检验数（含 $M$ 项）是最大的正检验数，算法优先级最高的就是把它们赶出基。

一旦人工变量全部离基（取值=0），$M\sum x_a=0$，目标函数恢复为原问题目标——此时得到的BFS就是原问题的BFS。

\textbf{如果赶不走怎么办？} 最优时若人工变量仍为正——说明在原问题的约束下，根本找不到让该人工变量为0的解。原问题不可行。

\textbf{方法（大M法标准流程）}

\begin{enumerate}
  \item \textbf{标准化后审视约束矩阵}：找出哪些行缺少 $+1$ 单位列。每个"松弛变量"列（$+1$）=天然的基候选；每个"剩余变量"列（$-1$）或缺变量行=需要人工变量。
  \item \textbf{引入人工变量}：在每个缺基的行添加人工变量 $x_a\ge 0$，系数 $+1$。人工变量只出现在自己的那行（系数列是单位向量）。
  \item \textbf{修改目标函数}：对 $\min$ 问题，目标加 $M\sum x_a$（$M$ 为充分大的正数）；对 $\max$ 问题，目标减 $M\sum x_a$。
  \item \textbf{建立初始单纯形表}：以人工变量+原有松弛变量为初始基。注意检验数行中人工变量的检验数含 $M$ 项——手工计算时将含 $M$ 和不含 $M$ 的部分分开写（通常写成两行或括号）。
  \item \textbf{正常执行单纯形法}：由于 $M$ 极大，含 $M$ 正系数的检验数最大——算法优先迭代赶走人工变量。
  \item \textbf{终止判断}：最优时检查人工变量状态——全为0则得原问题最优解；有人工变量 $>0$ 则原问题无可行解。
\end{enumerate}

\textbf{$M$ 在检验数中的作用：} 人工变量所在的行通过 $c_B$ 把 $M$ 传播到非基变量的检验数中，产生"驱动力"——含 $+M$ 系数的检验数最大，优先被选中进基，从而推动人工变量尽快离基。（课件 p5.30, p5.31）

\subsubsection{三种终止状态的判定}

\begin{enumerate}
  \item \textbf{得最优解}：所有检验数 $\le 0$ 且基变量中无人工变量
  \item \textbf{问题无界（无下界）}：某检验数 $>0$ 但该列所有系数 $\le 0$（没有正系数就没法算最小比值，意味着可以沿该方向无限走）
  \item \textbf{原问题无可行解}：所有检验数 $\le 0$ 但某个人工变量仍为正（基变量中还有人工变量 $>0$）——大M法的惩罚没把它赶走，说明原约束下该变量无法为0
\end{enumerate}

\subsubsection{两阶段法简介}

两阶段法是大M法的替代方案，避免 $M$ 带来的数值问题（实际计算中 $M$ 过大导致舍入误差）。
\begin{itemize}
  \item \textbf{第一阶段}：目标设为 $\min\sum x_a$（最小化人工变量之和），求解辅助LP。若最优值=0则所有人工变量可驱为0，得到原问题的一个初始BFS；若最优值 $>0$ 则原问题无可行解。
  \item \textbf{第二阶段}：去掉人工变量，从第一阶段得到的BFS出发，用原目标函数继续单纯形法求最优解。
\end{itemize}
两阶段法逻辑更清晰但步骤多；大M法一步到位但需处理 $M$ 符号。考试中大M法更常考。

\subsection{2.7 退化与Bland规则}

\subsubsection{退化是什么？}

\textbf{退化}指某个基变量取值为0（即 $x_{B_i}=0$）。退化的BFS中，正分量的个数严格小于 $m$。

\textbf{退化为什么麻烦？} 最小比值可能是 $0$（因为 $\bar{b}_i=0$ 且 $\bar{a}_{ik}>0$），导致步长为0——进基变量取代离基变量后，解不变（还是同一个点），但基变了。如果一连串0步长迭代后回到之前访问过的基——\emph{循环}发生，算法永不终止。

\textbf{Bland 规则（防循环）}

\textbf{进基}：在所有 $\sigma_j>0$ 的变量中，选下标最小的那个进基。

\textbf{离基}：若最小比值平局，选下标最小的基变量离基。

Bland规则是理论保障——严格证明有限步终止。实际计算中循环极其罕见，商业求解器通常不用Bland规则（有其他更高效的处理方式）。（课件 p5.9）

\subsection{本章速查}

\begin{itemize}
  \item \textbf{检验数} $\sigma_j = c_B^T B^{-1} a_j - c_j$（$\min$ 问题 $\sigma_j\le 0\;\forall j$ 为最优；$\sigma_k>0$ 且该列 $\le 0$ 则无界）
  \item \textbf{一次迭代}：进基（最大正检验数）$\to$ 离基（$\min\{b_i/a_{ik}\mid a_{ik}>0\}$）$\to$ 转轴（行变换使进基列为单位向量）
  \item \textbf{最小比值的几何意义}：沿进基方向走多远会碰到边界——取最近的那个边界
  \item \textbf{检验数行变换原理}：目标函数改写为 $z+\sum\sigma_j x_j = z_0$，与约束方程地位相同，行变换（Gauss消元）不改变解集，等价于对新基重算 $B'^{-1}A$
  \item \textbf{大M法因果链}：缺基 $\to$ 引入人工变量（凑单位阵）$\to$ 目标加 $M x_a$（惩罚驱赶）$\to$ $M$ 极大 $\to$ 含 $M$ 的检验数驱动算法优先赶走人工变量 $\to$ 人工变量全0后恢复原目标
  \item \textbf{两阶段法}：阶段一 $\min\sum x_a$ 找初始BFS，阶段二从该BFS求原问题最优
  \item \textbf{三种终止状态}：得最优解 / 问题无界 / 原问题无可行解（人工变量赶不走）
  \item \textbf{Bland规则}：进基选最小下标正检验数，离基平局选最小下标。防循环。
  \item \textbf{退化}：基变量=0导致步长=0，可能循环，Bland规则保证终止
  \item \textbf{影子价格}：最优表中 $c_B$ 列即为对偶变量（影子价格），可在单纯形表最优表中直接读出
\end{itemize}
